\documentclass[10pt,a4paper]{amsart}
\usepackage{amsmath,amssymb,amsthm,graphicx}
% For commutative diagrams we use Xy-pic
\usepackage[all]{xy}

% We use 2cell for 2-commutative diagrams.
\xyoption{2cell}
% Theorem environments.
%
\theoremstyle{plain}
\newtheorem{theorem}[subsection]{Theorem}
\newtheorem{proposition}[subsection]{Proposition}
\newtheorem{lemma}[subsection]{Lemma}

\theoremstyle{definition}
\newtheorem{definition}[subsection]{Definition}
\newtheorem{example}[subsection]{Example}
\newtheorem{exercise}[subsection]{Exercise}
\newtheorem{situation}[subsection]{Situation}

\theoremstyle{remark}
\newtheorem{remark}[subsection]{Remark}
\newtheorem{remarks}[subsection]{Remarks}

\numberwithin{equation}{subsection}

% Macros
%
\def\lim{\mathop{\mathrm{lim}}\nolimits}
\def\colim{\mathop{\mathrm{colim}}\nolimits}
\def\Spec{\mathop{\mathrm{Spec}}}
\def\Hom{\mathop{\mathrm{Hom}}\nolimits}
\def\Ext{\mathop{\mathrm{Ext}}\nolimits}
\def\SheafHom{\mathop{\mathcal{H}\!\mathit{om}}\nolimits}
\def\SheafExt{\mathop{\mathcal{E}\!\mathit{xt}}\nolimits}
\def\Sch{\mathit{Sch}}
\def\Mor{\mathop{\mathrm{Mor}}\nolimits}
\def\Ob{\mathop{\mathrm{Ob}}\nolimits}
\def\Sh{\mathop{\mathit{Sh}}\nolimits}
\def\NL{\mathop{N\!L}\nolimits}
\def\CH{\mathop{\mathrm{CH}}\nolimits}
\def\proetale{{pro\text{-}\acute{e}tale}}
\def\etale{{\acute{e}tale}}
\def\QCoh{\mathit{QCoh}}
\def\Ker{\mathop{\mathrm{Ker}}}
\def\Im{\mathop{\mathrm{Im}}}
\def\Coker{\mathop{\mathrm{Coker}}}
\def\Coim{\mathop{\mathrm{Coim}}}

% Boxtimes
%
\DeclareMathSymbol{\boxtimes}{\mathbin}{AMSa}{"02}

%
% Macros for moduli stacks/spaces
%
\def\QCohstack{\mathcal{QC}\!\mathit{oh}}
\def\Cohstack{\mathcal{C}\!\mathit{oh}}
\def\Spacesstack{\mathcal{S}\!\mathit{paces}}
\def\Quotfunctor{\mathrm{Quot}}
\def\Hilbfunctor{\mathrm{Hilb}}
\def\Curvesstack{\mathcal{C}\!\mathit{urves}}
\def\Polarizedstack{\mathcal{P}\!\mathit{olarized}}
\def\Complexesstack{\mathcal{C}\!\mathit{omplexes}}
% \Pic is the operator that assigns to X its picard group, usage \Pic(X)
% \Picardstack_{X/B} denotes the Picard stack of X over B
% \Picardfunctor_{X/B} denotes the Picard functor of X over B
\def\Pic{\mathop{\mathrm{Pic}}\nolimits}
\def\Picardstack{\mathcal{P}\!\mathit{ic}}
\def\Picardfunctor{\mathrm{Pic}}
\def\Deformationcategory{\mathcal{D}\!\mathit{ef}}
\usepackage{cleveref}
\crefname{originallemma}{original lemma}{original lemmas}
\Crefname{originallemma}{Original lemma}{Original lemmas}
\makeatletter
\newlabel{sites-modules-definition-morphism-locally-ringed-topoi}{{04HA}{source}{}{sourcepointer-1}{}}
\newlabel{sites-modules-definition-morphism-locally-ringed-topoi@cref}{{[originallemma][0][]04HA}{[1][0][]source}{}{}{}}
\makeatother
\crefname{originallemma}{original lemma}{original lemmas}
\Crefname{originallemma}{Original lemma}{Original lemmas}
\makeatletter
\newlabel{lemma-etale-site-locally-ringed}{{04HZ}{source}{}{sourcepointer-2}{}}
\newlabel{lemma-etale-site-locally-ringed@cref}{{[originallemma][0][]04HZ}{[1][0][]source}{}{}{}}
\makeatother
\crefname{originaltheorem}{original theorem}{original theorems}
\Crefname{originaltheorem}{Original theorem}{Original theorems}
\makeatletter
\newlabel{theorem-exactness-stalks}{{03PU}{source}{}{sourcepointer-3}{}}
\newlabel{theorem-exactness-stalks@cref}{{[originaltheorem][0][]03PU}{[1][0][]source}{}{}{}}
\makeatother
\crefname{originallemma}{original lemma}{original lemmas}
\Crefname{originallemma}{Original lemma}{Original lemmas}
\makeatletter
\newlabel{remarks-enough-points}{{040S}{source}{}{sourcepointer-4}{}}
\newlabel{remarks-enough-points@cref}{{[originallemma][0][]040S}{[1][0][]source}{}{}{}}
\makeatother
\crefname{originallemma}{original lemma}{original lemmas}
\Crefname{originallemma}{Original lemma}{Original lemmas}
\makeatletter
\newlabel{sites-modules-lemma-locally-ringed-morphism}{{04H9}{source}{}{sourcepointer-5}{}}
\newlabel{sites-modules-lemma-locally-ringed-morphism@cref}{{[originallemma][0][]04H9}{[1][0][]source}{}{}{}}
\makeatother
\crefname{originallemma}{original lemma}{original lemmas}
\Crefname{originallemma}{Original lemma}{Original lemmas}
\makeatletter
\newlabel{lemma-points-small-etale-site}{{04HU}{source}{}{sourcepointer-6}{}}
\newlabel{lemma-points-small-etale-site@cref}{{[originallemma][0][]04HU}{[1][0][]source}{}{}{}}
\makeatother
\crefname{originallemma}{original lemma}{original lemmas}
\Crefname{originallemma}{Original lemma}{Original lemmas}
\makeatletter
\newlabel{lemma-stalk-pullback}{{03Q1}{source}{}{sourcepointer-7}{}}
\newlabel{lemma-stalk-pullback@cref}{{[originallemma][0][]03Q1}{[1][0][]source}{}{}{}}
\makeatother
\title{The small morphism of ringed sites attached to a morphism of schemes is a morphism of locally ringed sites}
\author{The Stacks Project Authors}
\date{stacks-project:etale-cohomology.tex:local-snapshot-20261011; GNU FDL 1.2 or later}
\begin{document}
\maketitle
\noindent\textit{Editorial note.} All delivered mathematical text is the author's own, copied from the pinned Stacks Project chapter file etale-cohomology.tex; no statement or proof step was rewritten, reordered or removed.\par\medskip
\noindent\textit{Editorial note.} The proof is the author's complete verification of the defining condition; nothing is delegated.\par\medskip
\noindent\textit{Editorial note.} Every labelled reference occurring inside the counted spans is carried in assembly.external\_references as a non-load-bearing navigation pointer (same-chapter ones at their definition site, cross-chapter ones at the invocation inside the counted span), so no load-bearing mathematics is replaced by a stub. No bibliography citation occurs inside the counted spans.\par\medskip
\noindent\textit{Editorial note.} The checked chapter list for same-name collisions is the set of stacks chapters staged in this repository: chow, cohomology, divisors, etale-cohomology, flat, more-algebra, more-morphisms, morphisms, spaces-more-morphisms, varieties together with this card's chapter etale-cohomology.\par\medskip
\noindent\textit{Editorial note.} The preamble is the chapter's own native preamble excerpt already used for this source: the Xy-pic setup (native preamble lines 14-18) and the native theorem declarations with the mathematical macros (native preamble lines 154-217).\par\medskip
\noindent\textit{Allowed standard prerequisites (editorial).}\begin{itemize}
\item Morphism of ringed sites associated to a morphism of schemes. Same chapter, the morphism (f\_\{small\}, f\_\{small\}\textasciicircum{}sharp) of ringed sites attached to f : X -> Y. Supplies the object whose local-ringedness is verified.
\item Definition of a morphism of locally ringed sites and the condition on stalks. Modules on Sites, the definition of a morphism of locally ringed sites, which requires the induced maps on stalks at points to be local homomorphisms. Supplies the condition verified in the proof and the vocabulary of the statement.
\item Stalks of the structure sheaf at geometric points and the residue field description. Same chapter, the local ring O\textasciicircum{}\{sh\} at a geometric point and the identification of the stalk of the structure sheaf with it. Supplies the computation showing that the induced map on local rings is local.
\end{itemize}\medskip
\ifdefined\hypertarget\hypertarget{sourcepointer-1}{}\fi
\noindent\textit{Source reference.} Original lemma 04HA: sources/stacks-etale-cohomology/pinned-original-source.tex, source lines 5183--5183. Navigation pointer invoked at proof line 5183; the counted span contains the pointer only and the cited result is not reproduced inside it, so the pointer is registered as a non-load-bearing navigation stub in the established form for Stacks material. No bibliography citation is involved.\par\medskip
\ifdefined\hypertarget\hypertarget{sourcepointer-2}{}\fi
\noindent\textit{Source reference.} Original lemma 04HZ: sources/stacks-etale-cohomology/pinned-original-source.tex, source lines 4393--4393. Navigation pointer invoked at proof line 5191; the counted span contains the pointer only and the cited result is not reproduced inside it, so the pointer is registered as a non-load-bearing navigation stub in the established form for Stacks material. No bibliography citation is involved.\par\medskip
\ifdefined\hypertarget\hypertarget{sourcepointer-3}{}\fi
\noindent\textit{Source reference.} Original theorem 03PU: sources/stacks-etale-cohomology/pinned-original-source.tex, source lines 3546--3546. Navigation pointer invoked at proof line 5193; the counted span contains the pointer only and the cited result is not reproduced inside it, so the pointer is registered as a non-load-bearing navigation stub in the established form for Stacks material. No bibliography citation is involved.\par\medskip
\ifdefined\hypertarget\hypertarget{sourcepointer-4}{}\fi
\noindent\textit{Source reference.} Original lemma 040S: sources/stacks-etale-cohomology/pinned-original-source.tex, source lines 3583--3583. Navigation pointer invoked at proof line 5195; the counted span contains the pointer only and the cited result is not reproduced inside it, so the pointer is registered as a non-load-bearing navigation stub in the established form for Stacks material. No bibliography citation is involved.\par\medskip
\ifdefined\hypertarget\hypertarget{sourcepointer-5}{}\fi
\noindent\textit{Source reference.} Original lemma 04H9: sources/stacks-etale-cohomology/pinned-original-source.tex, source lines 5199--5199. Navigation pointer invoked at proof line 5199; the counted span contains the pointer only and the cited result is not reproduced inside it, so the pointer is registered as a non-load-bearing navigation stub in the established form for Stacks material. No bibliography citation is involved.\par\medskip
\ifdefined\hypertarget\hypertarget{sourcepointer-6}{}\fi
\noindent\textit{Source reference.} Original lemma 04HU: sources/stacks-etale-cohomology/pinned-original-source.tex, source lines 3610--3610. Navigation pointer invoked at proof line 5201; the counted span contains the pointer only and the cited result is not reproduced inside it, so the pointer is registered as a non-load-bearing navigation stub in the established form for Stacks material. No bibliography citation is involved.\par\medskip
\ifdefined\hypertarget\hypertarget{sourcepointer-7}{}\fi
\noindent\textit{Source reference.} Original lemma 03Q1: sources/stacks-etale-cohomology/pinned-original-source.tex, source lines 4649--4649. Navigation pointer invoked at proof line 5204; the counted span contains the pointer only and the cited result is not reproduced inside it, so the pointer is registered as a non-load-bearing navigation stub in the established form for Stacks material. No bibliography citation is involved.\par\medskip
\begin{lemma}
\label{lemma-morphism-locally-ringed}
Let $f : X \to Y$ be a morphism of schemes.
The morphism of ringed sites $(f_{small}, f_{small}^\sharp)$
associated to $f$ is a morphism of locally ringed sites, see
Modules on Sites,
Definition \ref{sites-modules-definition-morphism-locally-ringed-topoi}.
\end{lemma}

\begin{proof}
Note that the assertion makes sense since we have seen that
$(X_\etale, \mathcal{O}_{X_\etale})$ and
$(Y_\etale, \mathcal{O}_{Y_\etale})$
are locally ringed sites, see
Lemma \ref{lemma-etale-site-locally-ringed}.
Moreover, we know that $X_\etale$ has enough points, see
Theorem \ref{theorem-exactness-stalks}
and
Remarks \ref{remarks-enough-points}.
Hence it suffices to prove that $(f_{small}, f_{small}^\sharp)$
satisfies condition (3) of
Modules on Sites,
Lemma \ref{sites-modules-lemma-locally-ringed-morphism}.
To see this take a point $p$ of $X_\etale$. By
Lemma \ref{lemma-points-small-etale-site}
$p$ corresponds to a geometric point $\overline{x}$ of $X$.
By
Lemma \ref{lemma-stalk-pullback}
the point $q = f_{small} \circ p$ corresponds to the
geometric point $\overline{y} = f \circ \overline{x}$ of $Y$.
Hence the assertion we have to prove is that the induced map
of stalks
$$
(\mathcal{O}_Y)_{\overline{y}} \longrightarrow (\mathcal{O}_X)_{\overline{x}}
$$
is a local ring map. Suppose that $a \in (\mathcal{O}_Y)_{\overline{y}}$
is an element of the left hand side which maps to an element of the maximal
ideal of the right hand side. Suppose that $a$ is the equivalence class
of a triple $(V, \overline{v}, a)$ with $V \to Y$ \'etale,
$\overline{v} : \overline{x} \to V$ over $Y$, and $a \in \mathcal{O}(V)$.
It maps to the equivalence class of
$(X \times_Y V, \overline{x} \times \overline{v}, \text{pr}_V^\sharp(a))$
in the local ring $(\mathcal{O}_X)_{\overline{x}}$. But it is clear that
being in the maximal ideal means that pulling back $\text{pr}_V^\sharp(a)$
to an element of $\kappa(\overline{x})$ gives zero. Hence also pulling back
$a$ to $\kappa(\overline{x})$ is zero. Which means that $a$ lies in the
maximal ideal of $(\mathcal{O}_Y)_{\overline{y}}$.
\end{proof}

\end{document}
